\begin{tikzpicture}[font=\small,>=stealth,line join=round]
\draw[gray!45,rounded corners=3pt] (0,0) rectangle (14.2,9.4);
\node at (7.1,8.8) {不能代入不等于可以求导：分子并未趋于零};
\begin{scope}[shift={(3.5,4.4)},x=2cm,y=.6cm]
\draw[->,gray] (-1.5,0)--(1.65,0) node[above] {$t=x-1$};
\draw[->,gray] (0,-4.5)--(0,6.2) node[left] {$y$};
\draw[blue!75!black,very thick,samples=80,domain=-1.3:-.22] plot (\x,{1+1/\x});
\draw[blue!75!black,very thick,samples=80,domain=.22:1.3] plot (\x,{1+1/\x});
\draw[orange!85!black,thick,dashed] (-1.3,1)--(1.3,1);
\node[above left=8pt] at (0,1) {$1$};
\end{scope}
\node[anchor=west,align=left] at (8.8,5.2)
{蓝：$\dfrac{x}{x-1}=1+\dfrac1t$\\[18pt]
橙：误用后得到 $1$\\[18pt]
$t\to0^+:\ +\infty$\\[13pt]
$t\to0^-:\ -\infty$};
\node[text width=13.2cm,align=center] at (7.1,.65) {原式是 $1/0$ 型；分别求导得到的常数线，不能代表它的目标极限。};
\end{tikzpicture}
